\chapter{Implementation and Optimization}\label{chap:4}
{\itshape Implementing a real-time segmentation pipeline capable of processing over $120\text{ FPS}$ requires optimization tailored to low-level GPU hardware. The concrete C++ and CUDA implementation details are presented, detailing thread grid configurations, on-chip shared memory utilization, and hardware-accelerated math intrinsics. Additionally, CPU-GPU memory transfers are orchestrated through asynchronous streams to maximize pipeline throughput.}

\section{Hardware Environment and System Tuning}
To evaluate the pipeline fairly, all benchmarks were run on a standardized local hardware environment. The test system consists of an Intel Core Ultra 9 275HX CPU, 32 GB of system RAM, and an NVIDIA GeForce RTX 5070 Ti GPU with 12 GB of VRAM. Video input is acquired via a standard Chicony USB2.0 web camera capped at a hardware frame rate of 30 FPS, corresponding to an inter-frame capture latency of $33.3\text{ ms}$.

\subsection{System Tuning and Thermal Management}
To ensure deterministic latency and prevent frame-time variance, the hardware and operating system were tuned for sustained throughput. On the software side, the Windows scheduler was optimized via the registry (\texttt{Win32PrioritySeparation} set to \texttt{0x26}) to enable short, variable clock quantums, prioritizing foreground CPU slices for the camera input/output (I/O) and event loops \cite{microsoft_priority_separation_official, intel_thermal_13th_gen}. On the hardware side, rather than using max-frequency profiles that induce thermal throttling, a high performance profile was adopted to manage the thermal design power (TDP) between the CPU and GPU \cite{intel_thermal_13th_gen}, providing the stable thermal headroom necessary for continuous high-speed execution.

\subsection{Tooling and Code Quality}
The codebase is managed using the CMake build system \cite{cmake_official}, which handles the integration of C++20, OpenCV, GLFW, and the CUDA compiler (NVCC). To maintain clean and consistent code, \texttt{clang-format} is used for automated formatting and \texttt{clang-tidy} for static analysis \cite{clang_tools}. This setup helps prevent memory leaks and common C++ errors, ensuring benchmark results reflect the actual efficiency of the algorithms rather than implementation bugs.

\section{Threading Model and GPU Grid Configurations}
To execute 5D clustering in parallel on the GPU, the 2D image coordinates must be mapped to CUDA's thread and block grids \cite{CudaGuide}.

\subsection{2D Index Mapping}
For preprocessing, the \texttt{StridedDataPreprocessor} launches a 2D grid of thread blocks \cite{CudaGuide}. Each thread is mapped directly to an output pixel in the downsampled image \cite{cook2012cuda}. This bypasses the overhead of manual coordinate conversion and lets the hardware handle spatial locality automatically, as shown in Code Snippet~\ref{fig:thread_idx_code} \cite{CudaGuide}.

\begin{lstlisting}[caption={2D thread indexing and boundary check in \texttt{strided\_data\_\allowbreak preprocessor.cu} (lines 49--70).}, label={fig:thread_idx_code}]
// Map threads to output image coordinates
int out_x = blockIdx.x * blockDim.x + threadIdx.x;
int out_y = blockIdx.y * blockDim.y + threadIdx.y;

// Boundary check against sampled dimensions
if (out_x < out_cols && out_y < out_rows) {
    int x = out_x * stride;
    int y = out_y * stride;
    
    // Fetch raw BGR data and compute normalized coordinates
    int in_idx = y * cols + x;
    uchar3 bgr = frame_data[in_idx];
    
    float x_norm = static_cast<float>(x) * invCols;
    float y_norm = static_cast<float>(y) * invRows;
    // ... store in AoS memory layout ...
}
\end{lstlisting}

\subsection{Grid Heuristics for SM Occupancy}
A thread block size of $16 \times 16$ ($256$ threads per block) is selected. In the NVIDIA Blackwell architecture, this block size maximizes SM occupancy \cite{CudaGuide}. It provides a balanced distribution of registers per thread, allowing the hardware scheduler to hide memory latency by keeping warps highly active without running out of registers \cite{CudaGuide, cook2012cuda}.

\section{Implementation of Initialization Strategies}
The system supports two initialization strategies: Random and K-Means++.

\subsection{CPU Random Initialization}
For the baseline, the random initialization is implemented on the CPU. Because raw image data is staged in host memory via OpenCV \cite{Bradski2000}, selecting $K$ random starting pixels takes $O(K)$ time. Executing this on the CPU is near-instantaneous and avoids the PCIe transfer penalty of uploading raw index selections to the GPU, making it highly efficient.

\subsection{GPU-Accelerated K-Means++}
In contrast, the K-Means++ algorithm has a sequential complexity of $O(N \cdot K)$ \cite{arthur2007k}, which creates a major CPU bottleneck. This is offloaded to the GPU using custom CUDA kernels. 

For each of the $K$ seeds, a parallel distance-update kernel is executed. Each GPU thread calculates the squared distance between its assigned pixel and the newly selected centroid. A parallel prefix sum (reduction) is then executed in $O(\log N)$ using the Thrust library to build the probability distribution across the entire dataset \cite{bell2012thrust}. This allows for selecting the next seed in parallel, reducing K-Means++ initialization from several hundred milliseconds on the CPU to a sub-millisecond step on the GPU.

\section{Algorithmic Core (The Math Engines)}
The performance-critical path of the pipeline lies in the pixel-to-centroid assignment kernels.

\subsection{Templated Loop Unrolling}
To avoid code duplication between host and device, distance calculations utilize a templated math utility (\texttt{VectorMath<Dims>}). By passing the dimension count (5D) as a template parameter, the compiler is forced to perform complete loop unrolling (\texttt{\#pragma unroll}) \cite{CudaGuide}. The compiler expands the 5D feature loops into five sequential assembly operations, completely eliminating loop branch prediction overhead and instruction counters at runtime \cite{stroustrup2022tour, cook2012cuda}.

\subsection{Hardware-Accelerated Intrinsics}\label{subsect:intrinsics}
To achieve maximum frame rate, the active kernels use specialized GPU hardware features:
\begin{itemize}
    \item \textbf{Special Function Units (SFU):} In the \texttt{QuantumEngine}, features are mapped to cosine phase states. Instead of using the standard, high-precision C++ \texttt{cos()} function \cite{cppreference_cos}, the device intrinsic \texttt{\_\_cosf()} is utilized \cite{CudaGuide}. This compiles to direct hardware instructions on the GPU's SFUs, executing in a single clock cycle with negligible loss in precision \cite{CudaGuide}.
    \item \textbf{Fused Multiply-Add (FMA):} By enabling the \texttt{-use\_fast\_math} compiler flag, NVCC compiles distance calculations (such as $dist += diff \times diff$) into a single \texttt{\_\_fmaf()} instruction. This executes a multiply and an add in a single clock cycle, doubling floating-point throughput \cite{CudaGuide}.
\end{itemize}

The inner loop of the \texttt{QuantumEngine} assignment kernel demonstrates these hardware optimizations in Code Snippet~\ref{fig:quantum_sim_code}.

\begin{lstlisting}[caption={Quantum-inspired phase mapping and fast SFU cosine evaluation in \texttt{quantum\_engine.cu} (lines 68--82).}, label={fig:quantum_sim_code}]
// Map 5D features to phase orientations and compute global overlap
float target_prob = 1.0f;
#pragma unroll
for (int d = 0; d < 5; ++d) {
    float diff = (p[d] - s_centers[j * 5 + d]) * SCALE_FACTOR * PHASE_OFFSET;
    
    // Fast hardware-accelerated cosine intrinsic
    float cos_val = __cosf(diff);
    target_prob *= (cos_val * cos_val);
}
// Distance is the inverse of the overlap probability
float dist = 1.0f - target_prob;
\end{lstlisting}

\subsection{Scientific Justification for Single Precision (FP32)}
While high-precision scientific computing often requires double-precision (FP64), the pipeline is intentionally implemented using single-precision floats (FP32). This decision is justified by three hardware-level factors:
\begin{itemize}
    \item \textbf{Throughput:} Modern consumer GPUs have a $64:1$ ratio of FP32 to FP64 compute cores. Using FP32 allows the system to utilize the GPU's full parallel capacity \cite{CudaGuide, cook2012cuda}.
    \item \textbf{Precision Bounds:} The input pixel color channels are only 8-bit ($0-255$). The 24 bits of mantissa in an FP32 float are more than sufficient to represent these normalized values without convergence oscillations.
    \item \textbf{Memory Bandwidth:} FP32 halves the memory pressure compared to FP64, which is crucial when transferring and processing millions of 5D vectors per second \cite{CudaGuide, cook2012cuda}.
\end{itemize}

\section{Memory Optimization and Synchronization}
To achieve high-throughput clustering, the system minimizes global VRAM latency through a two-tier synchronization strategy.

\subsection{Two-Tier Shared Memory Atomics}
During the centroid update phase, thousands of threads must add their pixel features to a running sum for each cluster. If every thread executed an \texttt{atomicAdd} call directly on global VRAM, the memory controller would stall. Because standard atomic additions on global VRAM force the memory controllers to serialize requests, this contention generates hardware queues that destroy parallel throughput.

To prevent this, a two-tier aggregation kernel is implemented, shown in Code Snippet~\ref{fig:atomic_reduction_code}:
\begin{enumerate}
    \item \textbf{Local Reduction:} Threads within a block aggregate their pixel sums into low-latency, on-chip shared memory (\texttt{\_\_shared\_\_}) using fast atomic operations.
    \item \textbf{Global Reduction:} Once all threads in the block synchronize, a single designated thread per block executes an \texttt{atomicAdd} to write the block's sub-total into the global VRAM buffer. This reduces global memory traffic by a factor equal to the block size ($256\times$), eliminating VRAM contention \cite{CudaGuide, cook2012cuda}.
\end{enumerate}

\begin{lstlisting}[caption={Two-tier local-to-global atomic reduction in \texttt{base\_kmeans\_\allowbreak engine.cu} (lines 79--98).}, label={fig:atomic_reduction_code}]
// 1. Accumulate into block-local shared memory
atomicAdd(&s_counts[cluster], 1);
for (int d = 0; d < 5; ++d) {
    atomicAdd(&s_sums[cluster * 5 + d], samples[idx * 5 + d]);
}
__syncthreads();

// 2. Aggregate block-level results into global memory
if (threadIdx.x < k) {
    atomicAdd(&counts[threadIdx.x], s_counts[threadIdx.x]);
    for (int d = 0; d < 5; ++d) {
        atomicAdd(&newSums[threadIdx.x * 5 + d], s_sums[threadIdx.x * 5 + d]);
    }
}
\end{lstlisting}

\subsection{Cache Locality via Array of Structures (AoS)}
The 5D features are stored in an AoS memory layout: $[B_0, G_0, R_0, x_0, y_0, B_1, G_1, R_1,\allowbreak x_1, y_1...]$. While CUDA often favors a Structure of Arrays (SoA) for coalesced access \cite{CudaGuide, cook2012cuda}, the AoS layout was chosen to prioritize thread-level cache line utilization. 

Because the assignment kernels always access all five features of a pixel together, pulling a single feature channel automatically loads the adjacent color and spatial values into the same $L_1$/$L_2$ cache line. This layout guarantees a high cache hit rate for the remaining features and minimizes slow, independent global memory transactions, mitigating the latency of the strided global memory access \cite{CudaGuide}. Meaning that, for each pixel, only the first feature access potentially misses in cache; the remaining four features are served from the already-loaded cache line.

\section{The Visualization Pipeline: Asynchronous Staging}
To maintain a high-speed video feed, the visualization pipeline decouples GPU calculations from CPU rendering.

\subsection{Pinned Memory and DMA Transfers}
Host-side buffers are allocated using page-locked (pinned) memory via \texttt{cuda\allowbreak MallocHost}. This prevents the operating system from swapping these buffers to disk, allowing the GPU to use DMA to copy processed frames directly to system RAM without involving the CPU \cite{CudaGuide}. 

By combining pinned memory with non-blocking asynchronous streams (\texttt{cudaMemcpyAsync}), the visual segmented frame is transferred to the host while the CPU concurrently handles user interaction, camera I/O, and UI updates for the next frame. The stream is synchronized (\texttt{cudaStreamSynchronize}) only when the UI is ready to swap OpenGL textures, eliminating latency \cite{CudaGuide, cook2012cuda}.

\subsection{GPU-Native Classification and Colorization}
Rather than transferring raw integer labels back to the CPU and applying cluster colors there, the entire visualization pipeline is executed directly on the GPU. The custom \texttt{AssignPixelsKernel} performs classification and colorization in a single pass: each thread finds the nearest centroid for a pixel and immediately writes that cluster's assigned BGR color to the output display buffer. This reduces the visualization step to a single on-device write, bypassing expensive host-device data shuffles \cite{CudaGuide, cook2012cuda}.

\subsection{OpenGL Interoperability Rationale}
Pinned DMA staging was intentionally chosen over direct CUDA-OpenGL interop for two practical reasons:
\begin{itemize}
    \item \textbf{Thread Safety:} By using a pinned-host memory buffer as a staging area, the UI thread (OpenGL) and worker thread (CUDA) never access the same hardware resources simultaneously. This prevents race conditions without requiring complex thread synchronization locks.
    \item \textbf{API Portability:} Decoupling the GPU kernels from the graphics API ensures that the computational backend remains highly modular and can be easily ported to other windowing systems or headless servers without rewriting NVCC kernels.
\end{itemize}

\section{Centralized Parameter Management}
All mathematical boundaries and hardware constraints are controlled via a centralized configuration header (\texttt{constants.hpp}). This architectural choice ensures that the preprocessor, GPU engines, and UI always share the same parameters, as summarized in Table~\ref{tab:hyperparameters}.

\begin{table}[htbp]
\centering
\small
\caption[Critical System Hyperparameters and Functional Rationale]{Critical System Hyperparameters and Constraints}
\label{tab:hyperparameters}
    \begin{tabular}{lllp{6cm}}
    \hline
        \textbf{Category} & \textbf{Parameter} & \textbf{Value} & \textbf{Functional Rationale} \\ \hline
        Clustering & \texttt{FEATURE\_DIMS} & 5 & Defines the manifold as $[B, G, R, x, y]$. \\
         & \texttt{K\_MIN - K\_MAX} & 2 - 40 & Constrains model complexity for real-time stability. \\
         & \texttt{LEARN\_INTERVAL} & 15 & Performs "Lazy Learning" to reduce GPU thermal load. \\ \hline
        Normalization & \texttt{COLOR\_SCALE} & $1/255$ & Maps 8-bit integers to the $[0, 1]$ floating-point range. \\
         & \texttt{SPATIAL\_WEIGHT} & $10^{-6}$ & Regularizes $X/Y$ to prioritize color consistency. \\ \hline
        Quantum & \texttt{PHASE\_OFFSET} & 0.5 & Centers feature differences in the cosine phase space. \\
         & \texttt{SCALE\_FACTOR} & $\pi/2$ & Maps normalized distances to the $[0, \pi/2]$ phase range. \\ \hline
        Execution & \texttt{PROCESS\_WIDTH} & 640 px & Internal resolution to maximize kernel throughput. \\
          & \texttt{THREADS\_PER\_BLOCK} & 256 & Optimized block size for Blackwell SM occupancy. \\ \hline
        Metrics & \texttt{APPROX\_SUBSET\_SIZE} & 2,000 & Samples for Monte Carlo Silhouette estimation. \\ \hline
    \end{tabular}
\end{table}

\section{Performance-Aware Quality Metrics}
To evaluate segmentation quality, the system offloads three metrics (WCSS \cite{krzanowski1988criterion}, the Davies-Bouldin ($\mathcal{D}\mathcal{B}$) index \cite{davies1979cluster, xiao2017davies}, and the Silhouette Coefficient \cite{rousseeuw1987silhouettes, shutaywi2021silhouette}) to a dedicated background thread using the \texttt{BenchmarkRunner} class \cite{williams2019c++}. The optimized calculations are executed in \texttt{src/clustering/metrics.cpp} to process the telemetry data without blocking the real-time visualization pipeline.

Because calculating the exact Silhouette score for a $640 \times 480$ frame requires approximately $47.19$ billion calculations (Section~\ref{subsect:silhouette_complexity}), a Monte Carlo approximation using a random subset of $M=2,000$ pixels is implemented instead \cite{shutaywi2021silhouette}. This reduces processing time to under a millisecond, enabling real-time dashboard updates without user interface lag. These performance-aware quality metrics complete the implementation of the benchmarking pipeline. To ensure that these GPU-accelerated algorithms and asynchronous threads are mathematically correct and computationally stable, a comprehensive testing strategy is established, as detailed in the next chapter.
